\subsection{Topics 7 \& 8: Introduction to Networks and Application Layer – Exercise Sheet [Solutions]}
\subsection*{Q\#1: Consider an application that transmits data at a steady rate (for example, the sender generates an N-bit unit of data every k time unit, where k is small and fixed). Also, when such an application starts, it will continue running for a relatively long period of time. Answer the following questions, briefly justifying your answers:}
a) Would a packet-switched network or a circuit-switched network be more appropriate for this application? Why?
\newline
\textcolor{blue}{
A circuit-switched network would be well suited to the application because the application involves long sessions with predictable smooth bandwidth requirements. Since the transmission rate is known and not bursty, bandwidth can be reserved for each application session without significant waste. In addition, the overhead costs of setting up and tearing down connections are amortized over the lengthy duration of a typical application session.
}
\newline
b) Suppose that a packet-switched network is used and the only traffic in this network comes from such applications as described above. Furthermore, assume that the sum of the application data rates is less than the capacities of every link. Is some form of congestion control needed? Why?
\newline
\textcolor{blue}{
In the worst case, all the applications simultaneously transmit over one or more network links. However, since each link has sufficient bandwidth to handle the sum of all of the applications’ data rates, no congestion (very little queuing) will occur. Given such generous link capacities, the network does not need congestion control mechanisms.
}
\subsection*{Q\#2: Suppose there is exactly one switch between a sending host and a receiving host. The transmission rates between the sending host and the switch and between the switch and the receiving host are R1 = 10 Kbps and R2 = 2 Kbps, respectively. Assuming that the switch uses store-and-forward packet switching scheme, what is the total end-to-end delay to send a packet of length L = 1500 Bytes? (Ignore queuing, propagation, and processing delay).}
\textcolor{blue}{
At time t0 the sending host begins to transmit. At time t1 = L/R1 = (1500*8) / (10*1000), the sending host completes transmission and the entire packet is received at the router (no propagation delay). Because the router has the entire packet at time t1, it can begin to transmit the packet to the receiving host at time t1. At time t2 = t1 + (1500*8) / (2*1000), the router completes transmission, and the entire packet is received at the receiving host (again, no propagation delay). Thus, the end-to-end delay is L/R1 + L/R2 = 7.2 sec.
}
\subsection*{Q\#3: Suppose users share a 2Mbps link. Also, suppose each user transmits continuously at 1 Mbps when transmitting, but each user transmits only 20 percent of the time.}
a) When circuit-switching is used, how many users can be supported?
\newline
\textcolor{blue}{2 users can be supported because each user requires half of the link bandwidth.}
\newline
b) When packet-switching is used, why will there be essentially no queuing delay before the link if two or fewer users transmit at the same time? Why will there be a queuing delay if three users transmit at the same time?
\newline
\textcolor{blue}{Since each user requires 1Mbps when transmitting, if two or fewer users transmit simultaneously, a maximum of 2Mbps will be required. Since the available bandwidth of the shared link is 2Mbps, there will be no queuing delay before the link. Whereas, if three users transmit simultaneously, the bandwidth required will be 3Mbps which is more than the available bandwidth of the shared link. In this case, there will be a queuing delay before the link.}
\subsection*{Q\#4:}
a) How long does it take a packet of length L to propagate over a link of distance d, propagation speed s, and transmission rate R bps? Does this delay depend on packet length? Does this delay depend on the transmission rate?
$t = \frac{d}{s}$
\newline
No, it does not depend on the packet length.
\newline
No, it does not depend on the transmission rate.
\newline
b) How long does it take a packet of length 1,000 bytes to propagate over a link of distance 2,500 km, propagation speed $2.5 \times 10^8$ m/s, and transmission rate 2 Mbps?
\newline
\begin{center}
\(t = \frac{d}{s} = \frac{2500 \times 10^3}{2.5 \times 10^8} = 10 ms\)
\end{center}
\subsection*{Q\#5: Suppose Host A wants to send a large file to Host B. The path from Host A to Host B has three links of rates R1 = 500 kbps, R2 = 2 Mbps, and R3 = 1 Mbps.}
a) Assuming no other traffic in the network, what is the throughput for the file transfer?
\newline
\textcolor{blue}{
The throughput of the file transfer is equal to the bottleneck rate which is the smallest rate = 500 kbps.
}
\newline
b) Suppose the file is 4 million bytes. Given the above throughput, roughly how long will it take to transfer the file from Host A to Host B?
\newline
\textcolor{blue}{
500 kbps = 500,000 bps\
File size in bits = 4,000,000 * 8 = 32,000,000 bits\
Transfer time = $\frac{32000000}{500000} = 64$ seconds
}
\newline
c) Repeat (a) and (b), but now with R2 reduced to 100 kbps.
\newline
\textcolor{blue}{
Throughput = 100 kbps; Transfer time = $\frac{32000000}{100000} = 320$ seconds
}
\subsection*{Q\#6:}

List five non-proprietary Internet applications and the application-layer protocols that they use.
\textcolor{blue}{
\begin{enumerate}
    \item The Web: HTTP
    \item File transfer: FTP
    \item Remote login: Telnet
    \item E-mail: SMTP
    \item Domain Name System: DNS
\end{enumerate}
}
\subsection*{Q\#7: Assume you open a browser and enter \texttt{http://yourbusiness.com/about.html} in the address bar. What happens until the webpage is displayed? Provide details about the protocol(s) used and a high-level description of the messages exchanged.}
\textcolor{blue}{
In the first step, the domain name "yourbusiness.com" is translated into its IP address using DNS servers. This could be a multi-step process involving local DNS servers, root servers, TLD servers and the Authoritative DNS servers.
Once the translation is complete, the client contacts the HTTP server (using the IP address) for transferring the webpage. For the case of non-persistent HTTP connections, the steps are outlined below.
\begin{enumerate}
    \item The HTTP client process initiates a TCP connection to the server "yourbusiness.com" (using its IP) on port number 80, which is the default port number for HTTP. Associated with the TCP connection, there will be a socket at the client and a socket at the server.
    \item The HTTP client sends an HTTP request message to the server via its socket. The request message includes the path /about.html.
    \item The HTTP server process receives the request message via its socket, retrieves the object /about.html from its storage, encapsulates the object in an HTTP response message, and sends the response message to the client via its socket.
    \item The HTTP server process tells TCP to close the TCP connection.
    \item The HTTP client receives the response message. The TCP connection terminates. The message indicates that the encapsulated object is an HTML file. The client extracts the file from the response message, examines the HTML file, and displays the webpage on the client side.
\end{enumerate}
Application layer protocols: DNS and HTTP
\newline
Transport layer protocols: UDP for DNS; TCP for HTTP
}
\subsection*{Q\#8: What does a stateless protocol mean? Give a few examples. Is IMAP stateless? What about SMTP?}
\textcolor{blue}{
\begin{itemize}
  \item Stateless Protocols are the type of network protocols in which the Client sends a request to the server and the server responds according to the current state. It does not require the server to retain session information or status about each communicating partner for multiple requests. HTTP (Hypertext Transfer Protocol), UDP (User Datagram Protocol), DNS (Domain Name System) are examples of Stateless Protocols.
  \item IMAP is a stateful protocol because the IMAP server must maintain a folder hierarchy for each of its users, this state information persists across a particular user's successive accesses to the IMAP server. SMTP is a stateless protocol as the mail server does not maintain any connection with the client, it does not store any information about the client. If an email is asked to be sent twice, the server will resend it without saying that the email has been sent previously.
\end{itemize}
}
\subsection*{Q\#9: Suppose that your department has a local DNS server for all computers in the department. You are an ordinary user (i.e. not a network/system administrator). Can you determine if an external Web site was likely accessed from a computer in your department a couple of seconds ago? Explain.}
\textcolor{blue}{
Yes, we can use dig to query that Web site in the local DNS server. For example, “dig cnn.com” will return the query time for finding cnn.com. If cnn.com was just accessed a couple of seconds ago, an entry for cnn.com is cached in the local DNS cache, so the query time is 0 msec. Otherwise, the query time would be much larger.
}\\
\subsection*{Q\#10: What is the highest TCP/IP layer responsible for each of the following activities?}
a) Sending a frame to the next node.\\
\textcolor{blue}{
The Link Layer handles the transmission of frames between directly connected devices. It includes protocols like Ethernet and Wi-Fi.
}\\
b) Sending a packet from the source to the destination.\\
\textcolor{blue}{
The Network Layer is responsible for end-to-end packet delivery across networks. Protocols like IP (Internet Protocol) are used here.
}\\
c) Delivery of a long message from the source host to the destination host.\\
\textcolor{blue}{
The Transport Layer ensures complete and reliable message delivery, segmenting data into smaller chunks and reassembling them. Protocols include TCP (for reliability) and UDP.
}\\
d) Logging into a remote host using telnet or ssh protocol.\\
\textcolor{blue}{
The Application Layer provides network services to end users, including protocols like Telnet and SSH for remote login. It interfaces directly with software applications.
}